% Uniform all-gamma corollary, 22 October 2026.
\section{One polynomial for every scaling at least as strong as mean field}
\label{sec:uniform-supercritical}
The endpoint and spectral estimates control complementary scale ranges.
For the remaining bounded widths, a total-energy bound suffices. Placing
all three ranges in one auxiliary function class produces one sample
event for every width and scale. Parameterize the raw
Gaussian prior by $0<t\le1$, with coordinate variances
\[
 \frac{t^2}{d}\quad\hbox{in the first layer},\qquad
 \frac{t^2}{N}\quad\hbox{in internal layers},\qquad
 \frac{t^2}{N^2}\quad\hbox{at the readout}.
\]
The scaling exponent $\gamma\ge1$ corresponds to
$t=N^{(1-\gamma)/2}$. Allowing every $t\in(0,1]$ also permits the scale
to be chosen independently at each width. Write
$\nu_{N,t}=\mu_{N,\sigma=t/\sqrt N}$ for this raw law.

\begin{lemma}[Uniformity of the compression ingredients]
\label{lem:uniform-scale-ingredients}
Let $A_f=A$ be the constant in
Proposition~\ref{prop:unbounded-target-mass}. At every $N\ge C$ and
$0<t\le1$,
\[
 \nu_{N,t}(U_{N,f})\ge e^{-A_fN/(2t^2)}.
\]
The mean-field functional compression proposition remains true with the
same polynomials $P_{L,\lambda},J_{L,\lambda}$ and raw exceptional bound
$(2H+1)e^{-NK}$ whenever $N\ge P_{L,\lambda}(d,K)$, uniformly for
every $t\in(0,1]$.
Also, whenever
\[
 t\le a_K/8,\qquad t^{-2}\ge8(d+1),
\]
the spectral functional compression at energy $(L+3)K$ holds at every
width, with raw exceptional mass at most $Le^{-NK/(8t^2)}$.
\end{lemma}
\begin{proof}
The target-cylinder proof uses scalar Gaussian variance $t^2$ in its
constrained coordinates and $t\le1$ in its deterministic error bound;
replacing $N^\gamma$ by $N/t^2$ gives the first assertion directly.

In the endpoint argument the conditional preactivation covariance is
now $t^2Q_{s-1}$ and the first one is $t^2Q_0$. Apply the same one-step
construction to these covariances. Their diagonal budgets do not increase.
If $Q_{s-1}'$ approximates $Q_{s-1}$, then
\[
 \|t^2Q_{s-1}'-t^2Q_{s-1}\|_{\max}
 \le\|Q_{s-1}'-Q_{s-1}\|_{\max}.
\]
The weighted feature coefficients for the next preactivation are
multiplied by $t^2$. Thus Gaussian incoming weights are multiplied by
$t$, atom incoming weights by $t^2$, and all existing upper bounds on
weights and feature magnitudes remain valid. Initially the input-feature
coefficients are $t^2/d$. The readout energy is
$E_v=t^2\chi_N^2/N$, whose upper tail is no larger than at $t=1$.
Hence the same polynomial thresholds, compressed-energy bound, and
raw exceptional probability apply. Although the reconstructed networks
may depend on $t$, the single auxiliary energy class containing them does
not depend on it.

For spectral compression, the weighted first-layer energy and readout
energy are $(t^2/N)\chi^2_{N(d+1)}$ and $(t^2/N)\chi^2_N$;
the internal entries have variance $t^2/N$. The two stated inequalities
are exactly the hypotheses used in the spectral proof. Its singular-value,
RMS-activation, and readout estimates therefore give the displayed tail,
without a lower bound on $N$.
\end{proof}

Choose the following budgets, all independent of learner width and scale:
\begin{align*}
 K_s&=4A_f+8\log(2L/\varepsilon),&
 t_*&=\min\{a_{K_s}/8,[8(d+1)]^{-1/2}\},\\
 K_m&=2+A_f/t_*^2+\log\{2(2H+1)/\varepsilon\},&
 N_*&=\left\lceil\max\{1,P_{L,\lambda}(d,K_m)\}\right\rceil,\\
 D_*&=d+2+(L-2)N_*,&
 J_0&=2A_f+2D_*\log2+4\log(2/\varepsilon),\\
 J_*&=(L+3)K_s+J_{L,\lambda}(d,K_m)+J_0.
\end{align*}
Each budget, including $t_*^{-1}$, is polynomially bounded in
$d,C,\log(1/\varepsilon)$ for fixed $L,\lambda$.

\begin{theorem}[All widths and all scales at least as strong as mean field]
\label{thm:uniform-all-scales}
Fix a depth $L\ge2$ and an activation $\phi$ with $\phi(0)=0$ and
Lipschitz constant at most $\lambda$. Suppose a Boolean target has a
width-$C$, depth-$L$ implementation with weights bounded by $C$ and
strict global margin $1/2$. Let the original fixed-$\phi$ Gaussian prior
have the variances displayed above, condition it on an absolute margin
$1/2$ at every Boolean input, and then condition on exact training labels.
For every input distribution on the Boolean cube and every
$0<\varepsilon,\delta<1$, one event of training-sample probability at
least $1-\delta$ gives posterior average sign error at most $\varepsilon$
\emph{simultaneously for every integer $N\ge C$ and every $0<t\le1$},
provided
\[
 n\ge\frac2\varepsilon
 \left[\mathcal B(J_*)+\log\frac1\delta\right].
\]
In particular, one fixed-depth polynomial sample bound works for every
$\gamma\ge1$, including $\gamma=1$, without a polynomial degree or
constant that deteriorates as $\gamma\downarrow1$.
ReLU, exact GELU, SiLU, tanh, and every other fixed Lipschitz activation
vanishing at zero are covered whenever the stipulated teacher exists.
For the correct label to have posterior probability strictly greater
than $2/3$ outside a test set of mass at most $\varepsilon$, apply the
bound with average-error accuracy $\varepsilon/3$.
\end{theorem}
\begin{proof}
Let $\mathcal F_*$ be the single auxiliary strict-margin function class
of energy at most $J_*$. We show that at every $N\ge C$ and $0<t\le1$,
the raw probability of margin-separated functions outside this class,
divided by the raw target probability, is at most $\varepsilon/2$.

If $t\le t_*$, spectral compression gives the ratio bound
\[
 L\exp\left[-\frac N{t^2}(K_s/8-A_f/2)\right]
 \le\varepsilon/2.
\]
This applies at every width.
If $t\ge t_*$ and $N\ge N_*$, the uniform endpoint compression gives
\[
 (2H+1)\exp\left[-N\left(K_m-\frac{A_f}{2t^2}\right)\right]
 \le\varepsilon/2.
\]

For the remaining bounded widths $C\le N<N_*$, use total energy instead
of a union over the continuum of possible scales. The parameter count is
$p=N(d+2)+(L-2)N^2$, so $p/N\le D_*$. The normalized energy has law
$\kappa_N=(t^2/N)\chi_p^2$. Its exponential moment at $N/(4t^2)$ is
$2^{p/2}$, and therefore
\begin{align*}
 \frac{\nu_{N,t}\{\kappa_N>J_0\}}{\nu_{N,t}(U_{N,f})}
 &\le \exp\left\{-\frac N{t^2}(J_0/4-A_f/2)
                                      +\frac p2\log2\right\}\\
 &\le \exp\left\{-\frac N{t^2}
                  (J_0/4-A_f/2-D_*\log2/2)\right\}
 \le\varepsilon/2.
\end{align*}
Here $t\le1$ and $N\ge1$ justify the last two inequalities. A globally
margin-separated draw with $\kappa_N\le J_0$ itself belongs to the
enlarged auxiliary class. Thus all three regimes are contained in the
same $\mathcal F_*$ apart from the asserted relative mass.

The target must belong to $\mathcal F_*$, since otherwise the relative
mass would be at least one. Normalization by the global margin event
cancels, and conditioning on any noiseless training labels preserves the
upper bound $\varepsilon/2$ on posterior mass outside $\mathcal F_*$.
The capacity theorem bounds $\log|\mathcal F_*|$ by
$\mathcal B(J_*)$. A finite-class union bound ensures, on one sample
event of probability at least $1-\delta$, that every consistent member
has true error at most $\varepsilon/2$. That event refers only to the
single function class, so it works for the entire continuum of scales
and every learner width simultaneously. Adding the posterior mass of
the complement gives the claimed average-error bound.
\end{proof}
